% Generated by scripts/assemble.py; corrections belong in source/corrections.json.
\examhead{155}{华南理工大学 2024 年数学分析试卷}{147}
% p148-b19; source PDF page 148
\par\Needspace{6\baselineskip}
\blockanchor{p148-b19}
\question{1}(13 分) 已知 $a, b > 0$, 且 $c = \lim_{x \to 0} \left( \frac{4}{\ln(1 + x)} - \frac{4}{x} \right)$, 定义函数\par
\nopagebreak[4]
% p148-b20; source PDF page 148
\blockanchor{p148-b20}
\displayformula{f (x) = \left\{ \begin{array}{l l} \frac {\sin (a x)}{x}, & x <   0; \\ c, & x = 0; \\ (1 + b x) ^ {\frac {1}{a}}, & x > 0. \end{array} \right.}
% p148-b21; source PDF page 148
\blockanchor{p148-b21}
若 $f(x)$ 在 $(-\infty, +\infty)$ 内处处连续，求常数 $a, b$ 和 $c$.\par
% p149-b00; source PDF page 149
\par\Needspace{6\baselineskip}
\blockanchor{p149-b00}
\question{2}(13 分) 设 $f(x)$ 在 $[a, b]$ 上连续, 在 $(a, b)$ 内可导. 证明: 存在 $\xi \in (a, b)$, 使得\par
\nopagebreak[4]
% p149-b01; source PDF page 149
\blockanchor{p149-b01}
\displayformula{f (\xi) + f ^ {\prime} (\xi) = \frac {e ^ {b} f (b) - e ^ {a} f (a)}{e ^ {b} - e ^ {a}}.}
% p149-b02; source PDF page 149
\par\Needspace{6\baselineskip}
\blockanchor{p149-b02}
\question{3}(13 分) 证明函数\par
\nopagebreak[4]
% p149-b03; source PDF page 149
\blockanchor{p149-b03}
\displayformula{f (x, y) = \left\{ \begin{array}{l l} x y \cos \frac {1}{\sqrt {x ^ {2} + y ^ {2}}}, & (x, y) \neq (0, 0); \\ 0, & (x, y) = (0, 0). \end{array} \right.}
% p149-b04; source PDF page 149
\blockanchor{p149-b04}
在点 $(0,0)$ 连续且偏导数存在, 但是偏导数在点 $(0,0)$ 不连续, 而 $f(x,y)$ 在点 $(0,0)$ 可微.\par
% p149-b05; source PDF page 149
\blockanchor{p149-b05}
\question{4}(13 分) 将方程 $\frac{\partial^2 z}{\partial x^2} + \frac{\partial^2 z}{\partial y^2} = 0$ 变换为极坐标形式.\par
% p149-b06; source PDF page 149
\par\Needspace{6\baselineskip}
\blockanchor{p149-b06}
\question{5}(14 分) 计算曲线积分 $\int_{L} y \, \mathrm{d}x + z \, \mathrm{d}y + x \, \mathrm{d}z$, 其中 $L$ 是曲线\par
\nopagebreak[4]
% p149-b07; source PDF page 149
\blockanchor{p149-b07}
\displayformula{\frac {x ^ {2}}{a ^ {2}} + \frac {y ^ {2}}{b ^ {2}} + \frac {z ^ {2}}{c ^ {2}} = 1, \frac {x}{a} + \frac {z}{c} = 1, x \geq 0, y \geq 0, z \geq 0.}
% p149-b08; source PDF page 149
\blockanchor{p149-b08}
$(a>0,b>0,c>0$  为常数 $)$ 从点 $(a,0,0)$  到 $(0,0,c)$ .\par
% p149-b09; source PDF page 149
\blockanchor{p149-b09}
\question{6}(14 分) 证明黎曼 zeta 函数 $\zeta(x) = \sum_{n=1}^{\infty} \frac{1}{n^x}$ 在 $(1, +\infty)$ 上连续且无穷次可微.\par
% p149-b10; source PDF page 149
\blockanchor{p149-b10}
\question{7}(14 分) 解答如下问题:\par
% p149-b11; source PDF page 149
\blockanchor{p149-b11}
(1) 用定义证明 $\lim_{n\to \infty}\sin n$ 不存在.\par
% p149-b12; source PDF page 149
\par\Needspace{6\baselineskip}
\blockanchor{p149-b12}
(2) 证明数列 $\{x_{n}\}$ 极限存在, 其中\par
\nopagebreak[4]
% p149-b13; source PDF page 149
\blockanchor{p149-b13}
\displayformula{x _ {n} = \left(1 - \frac {1}{2 ^ {2}}\right) \left(1 - \frac {1}{3 ^ {2}}\right) \dots \left(1 - \frac {1}{(n + 1) ^ {2}}\right) = \prod_ {k = 1} ^ {n} \left(1 - \frac {1}{(k + 1) ^ {2}}\right), n = 1, 2, \dots .}
% p149-b14; source PDF page 149
\blockanchor{p149-b14}
并求极限 $\lim_{n\to \infty}x_n$\par
% p149-b15; source PDF page 149
\par\Needspace{6\baselineskip}
\blockanchor{p149-b15}
\question{8}(14 分) 设函数 $f(x)$ 在 $(x_0 - r, x_0 + r)$ 上无穷次可导(其中 $r \in (0, +\infty)$), 且存在常数 $M > 0$, 使得\par
\nopagebreak[4]
% p149-b16; source PDF page 149
\blockanchor{p149-b16}
\displayformula{| f ^ {(n)} (x) | \leq M \frac {n !}{r ^ {n}}, \forall x \in (x _ {0} - r, x _ {0} + r), \forall n = 0, 1, 2, \dots}
% p149-b17; source PDF page 149
\blockanchor{p149-b17}
(1) 证明 $f(x)$ 在 $x_0$ 点的 Taylor 级数 $\sum_{n=0}^{\infty} \frac{f^{(n)}(x_0)}{n!}(x - x_0)^n$ 在 $(x_0 - r, x_0 + r)$ 上逐点收敛于 $f(x)$.\par
% p149-b18; source PDF page 149
\blockanchor{p149-b18}
(2) 证明 $\sum_{n=0}^{\infty} \frac{f^{(n)}(x_0)}{n!}(x - x_0)^n$ 在 $(x_0 - r, x_0 + r)$ 上内闭一致收敛于 $f(x)$.\par
% p149-b19; source PDF page 149
\par\Needspace{6\baselineskip}
\blockanchor{p149-b19}
\question{9}(14 分) 设 $f(x)$ 在 $[a, b]$ 上黎曼可积, 证明黎曼引理\par
\nopagebreak[4]
% p149-b20; source PDF page 149
\blockanchor{p149-b20}
\displayformula{\lim _ {\lambda \rightarrow + \infty} \int_ {a} ^ {b} f (x) \cos (\lambda x) \mathrm{d} x = 0.}
% p149-b21; source PDF page 149
\blockanchor{p149-b21}
\question{10}(14 分) 证明 $\int_0^{+\infty} \frac{\sin xy}{x} \, \mathrm{d}x$ 关于 $y$ 在 $[y_0, +\infty)$ ($y_0 > 0$) 上一致收敛, 但在 $(0, +\infty)$ 上非一致收敛.\par
% p149-b22; source PDF page 149
\blockanchor{p149-b22}
\question{11}(14 分) 设 $f(x)$ 在 $[a, b]$ 上分段连续, 即存在 $[a, b]$ 的一个有限分割 $a = x_0 < x_1 < x_2 < \cdots < x_n = b$ (其中 $n$ 为固定整数), 使得 $f$ 在每个区间 $(x_{i-1}, x_i)$ 上连续且分点 $x_i$ 处都存在左右极限. 证明 $f(x)$ 在 $[a, b]$ 上黎曼可积.\par
